\par\noindent\begin{minipage}{\linewidth}
\begin{DefinitionBox}{Polynomfunktion und Grad}{real_functions_polynomial}
  Eine \emph{Polynomfunktion} \EnglishTerm{polynomial function} besitzt
  für $x\in\R$ einen Term der Form
  \[
    p(x)=a_nx^n+a_{n-1}x^{n-1}+\dots+a_1x+a_0,
    \qquad a_0,\dots,a_n\in\R.
  \]
  Für $a_n\ne0$ heißt $n\in\N_0$ der \emph{Grad} \EnglishTerm{degree}
  und $a_n$ der \emph{Leitkoeffizient} \EnglishTerm{leading coefficient}.
  Das Nullpolynom $p(x)=0$ erhält hier keinen Grad.
\end{DefinitionBox}
\begin{ExampleBox}[before skip=0.3\baselineskip]{Den Term lesen}{real_functions_polynomial_reading}
  Für $p(x)=2x^4-3x^2+1$ lassen sich Grad, Achsenabschnitt und Symmetrie
  unmittelbar erkennen:
  \[
    \deg p=\Cloze{1cm}{4},\qquad
    p(0)=\Cloze{1cm}{1},\qquad
    p(-x)=\Cloze{2cm}{p(x)}.
  \]
  Fehlende Potenzen haben den Koeffizienten null. Nur gerade Potenzen
  einschließlich der Konstanten ergeben eine gerade Funktion.
\end{ExampleBox}
\end{minipage}


\subsection{Vorgegebene Punkte durch ein Polynom verbinden}
\label{subsec:real_functions_interpolation}
Bei der \emph{Interpolation} \EnglishTerm{interpolation} soll der Graph
vorgegebene Punkte exakt durchlaufen. Die unbekannten Koeffizienten
lassen sich mit einem linearen Gleichungssystem bestimmen.
\begin{ExampleBox}{Parabel durch drei Punkte}{real_functions_interpolation_example}
  Gesucht ist ein Polynom vom Grad höchstens zwei durch
  $P=(0,0)$, $Q=(1,1)$ und $R=(2,0)$.
  \par\smallskip
  \begin{minipage}[t]{0.52\linewidth}
    \vspace{0pt}
    Wir wählen den Ansatz $p(x)=ax^2+bx+c$ und setzen die Punkte ein.
    \begin{ClozeArea}[17]{9}
      \[
        \begin{aligned}
          p(0)=0&\Rightarrow c=0,\\
          p(1)=1&\Rightarrow a+b=1,\\
          p(2)=0&\Rightarrow4a+2b=0.
        \end{aligned}
      \]
      Daraus folgen $a=-1$, $b=2$ und somit
      \[
        p(x)=-x^2+2x.
      \]
    \end{ClozeArea}
  \end{minipage}\hfill
  \begin{minipage}[t]{0.44\linewidth}
    \vspace{0pt}\centering
    \CachedTikzFigure{content/resources/chapters/real_functions/tikz/functions_interpolation.tex}
  \end{minipage}
\end{ExampleBox}
\begin{TheoremBox}{Eindeutigkeit der Polynominterpolation}{real_functions_interpolation_unique}
  Zu $n+1$ Punkten $(x_0,y_0),\dots,(x_n,y_n)$ mit paarweise verschiedenen
  $x$-Werten gibt es genau ein Polynom vom Grad \emph{höchstens} $n$,
  das alle Punkte durchläuft.
\end{TheoremBox}
